\begin{equation*}
\max \left\{m (f _ {1}, f _ {2}, \dots , f _ {n}) \right\} = \lambda_ {n + 1},\tag{4.38}
\end{equation*}

这就是所要证明的。

现在我们可着手证明关于埃尔米特算子特征值的一个有趣而且也有用的定理。考虑一界限任意容积 $R'$ 的曲面 $S, R'$ 本身位于一大平行六面体 $B$ 内(图23)。设算子 $H = -\triangle$ 作用于定义在容积 $R'$ 内的函数，此外并假定对于属于容许函数类的波动函数 $\psi$ 在 $S$

\begin{figure}[htbp]
\centering
\includegraphics[width=0.52\textwidth]{fig23.pdf}
\par\vspace{2pt}
{\small 图 23.}
\end{figure}

上的边界条件有 $\psi = 0$ 的形式，那么在容积 $R'$ 内

\begin{equation*}
H \psi = \lambda^ {\prime} \psi .\tag{4.39}
\end{equation*}

类似地，在 $B$ 的内部设 $H = -\triangle$ ，而在边界 $B$ 上 $\varphi = 0$ 。于是

\begin{equation*}
H \varphi = \lambda \varphi .\tag{4.40}
\end{equation*}

现将关于(4.39)和(4.40)的特征值集合排列成

\begin{equation*}
\begin{array}{l} \lambda_ {1} ^ {\prime} \leqslant \lambda_ {2} ^ {\prime} \leqslant \lambda_ {3} ^ {\prime} \leqslant \dots , \\ \lambda_ {1} \leqslant \lambda_ {2} \leqslant \lambda_ {3} \leqslant \dots . \end{array}
\end{equation*}

则有下面的定理。

\textbf{定理 4.5}\quad 对所有 $n$ ，下面的不等式成立：

\begin{equation*}
\lambda_ {n} ^ {\prime} \geqslant \lambda_ {n}.\tag{4.41}
\end{equation*}

\textbf{证明}\quad 用 $\varphi'$ 记下面的函数：

\begin{equation*}
\varphi^ {\prime} = \left\{ \begin{array}{l l} \psi & \text {在} R ^ {\prime} \text {内}, \\ 0 & \text {在} S \text {和} B \text {之间\text{。}} \end{array} \right.\tag{4.42}
\end{equation*}

其次我们注意

\begin{equation*}
\lambda \left(\varphi^ {\prime}\right) = - \frac {\int_ {B} \left(\varphi^ {\prime}\right) ^ {*} \triangle \varphi^ {\prime} d \tau}{\int_ {B} \left(\varphi^ {\prime}\right) ^ {*} \varphi^ {\prime} d \tau} = \frac {\int_ {B} \left| \nabla \varphi^ {\prime} \right| ^ {2} d \tau}{\int_ {B} \left(\varphi^ {\prime}\right) ^ {*} \varphi^ {\prime} d \tau},\tag{4.43}
\end{equation*}

最后一个等式是通过分部积分而得到的，其中的一项由于所加的边界条件而消失了。类似地

\begin{equation*}
\lambda^ {\prime} (\psi) = \frac {\int_ {B} | \nabla \psi | ^ {2} d \tau}{\int_ {B} \psi^ {*} \psi d \tau}.\tag{4.44}
\end{equation*}

因在 $R$ 内 $\nabla \psi = \nabla \varphi'$ ，又因在 $S$ 外两函数等于零，由(4.43)和(4.44)得

\begin{equation*}
\lambda^ {\prime} (\psi) = \lambda \left(\varphi^ {\prime}\right).\tag{4.45}
\end{equation*}

由所有函数 $\varphi$ 组成的函数类，用 $G$ 表示；而由所有函数 $\varphi'$ 组成的函数类用 $G'$ 表示。正如已证明的那样，当 $\varphi \in G$ 时，

\begin{equation*}
\min \lambda (\varphi) = \lambda_ {1};
\end{equation*}

而当 $\varphi' \in G'$ 时，

\begin{equation*}
\min \lambda (\varphi ^ {\prime}) = \lambda_ {1} ^ {\prime}.
\end{equation*}

因为 $G' \subset G(G'$ 是 $G$ 的子集)，故 $G$ 比 $G'$ 含有更多的函数。因而

\begin{equation*}
\min \lambda (\varphi) = \lambda_ {1} \leqslant \min \lambda \left(\varphi^ {\prime}\right) = \lambda_ {1} ^ {\prime}.\tag{4.46}
\end{equation*}

我们的目的是要证明对 $\lambda_{2}$ 和 $\lambda_{2}^{\prime}$ 有类似的关系。用

\begin{equation*}
\lambda_ {2} = \min \lambda (\varphi)
\end{equation*}

记对所有与 $\varphi_{1}$ 正交且属于 $G$ 的 $\varphi$ 的 $\lambda$ 之极小值；而用

\begin{equation*}
\lambda_ {2} ^ {\prime} = \min \lambda (\varphi^ {\prime})
\end{equation*}

记对所有与 $\varphi_{1}^{\prime}$ 正交且属于 $G^{\prime}$ 的 $\varphi^{\prime}$ 的 $\lambda$ 之极小值。由此还不能得出 $\lambda_{2}^{\prime} \geqslant \lambda_{2}$ 的结论，因即使 $G^{\prime} \subset G$ ，但由正交于 $\varphi_{1}$ 的函数 $\varphi$ 组成的 $G$ 的子集可以小于由正交于 $\varphi_{1}^{\prime}$ 的函数 $\varphi^{\prime}$ 组成的 $G^{\prime}$ 的子集。可是定理4.4在这里是有用的。

任取函数 $f_{1}$ 。假定 $m^{\prime}(f_{1})$ 是对属于 $G^{\prime}$ 且正交于 $f_{1}$ 的所有 $\varphi^{\prime}$ 的 $\lambda (\varphi^{\prime})$ 之极小值；而 $m(f_{1})$ 是对属于 $G$ 且正交于 $f_{1}$ 的所有 $\varphi$ 的 $\lambda (\varphi)$ 之极小值，即

\begin{equation*}
\begin{array}{r l} m ^ {\prime} (f _ {1}) & = \min \lambda (\varphi^ {\prime}), \\ m (f _ {1}) & = \min \lambda (\varphi). \end{array}\tag{4.46a}
\end{equation*}

现在可断言
\begin{equation*}
m ^ {\prime} (f _ {1}) \geqslant m (f _ {1}),\tag{4.47}
\end{equation*}

因为所有 $\varphi$ 和 $\varphi'$ 都正交于同一个函数，故正交于 $f_{1}$ 的所有 $\varphi$ 的集合包含正交于 $f_{1}$ 的所有 $\varphi'$ 为其子集。

现在应用定理 4.4。由不等式(4.47)得

\begin{equation*}
\max \left[ m ^ {\prime} \left(f _ {1}\right) \right] = \lambda_ {2} ^ {\prime} \geqslant \max \left[ m \left(f _ {1}\right) \right] = \lambda_ {2},
\end{equation*}

即
\begin{equation*}
\lambda_ {2} ^ {\prime} \geqslant \lambda_ {2}.
\end{equation*}

现证明
\begin{equation*}
\lambda_ {3} ^ {\prime} \geqslant \lambda_ {3}.
\end{equation*}

为此，定义两个函数 $f_{1}$ 和 $f_{2}$ ，使得对属于 $G'$ 且正交于 $f_{1}$ 和 $f_{2}$ 的所有 $\varphi'$ ， $m'(f_1, f_2) = \min \lambda(\varphi')$ ；而对属于 $G$ 且正交于 $f_{1}$ 和 $f_{2}$ 的所有 $\varphi$ ， $m(f_1, f_2) = \min \lambda(\varphi)$ 。

类似于前面的讨论可得

\begin{equation*}
\lambda_ {3} ^ {\prime} \geqslant \lambda_ {3}.
\end{equation*}

重复这个方法，我们就能对所有 $n$ 建立一般的关系式

\begin{equation*}
\lambda_ {n} ^ {\prime} \geqslant \lambda_ {n}.
\end{equation*}

这样就完成了定理的证明。

由这个定理可推出两个十分重要的推论：

1) 由对于平行六面体问题的解答知 $\varphi$ 为完备正交系。这是因为 $\lambda_{1} < \infty$ ，而 $\lim_{n \to \infty} \lambda_{n} = \infty$ 。因此， $\lim_{n \to \infty} \lambda_{n}' = \infty$ （由定理4.5），同时 $\lambda_{1}'$ 有限，故 $\psi$ 也成为完备正交系，即埃尔米特算子 $H = -\triangle$ 满足边界条件 $\psi = 0$ 的特征函数构成完备正交系。

2) 考虑有确定特征函数的膜片，如果对它加以补充关系，例如在确定点上加上不大的负荷(即附加力)，则所有特征值都增加了。

一般地，增加问题中关系的数目，就会引起特征值的增加，这也适用于算子 $H = -\triangle$ 。

\section{正算子}

继续研究希尔伯特空间。考虑一个特殊类型的算子，即所谓正算子。

如果任一算子 $H$ ，对任意函数 $\varphi$ 有

\begin{equation*}
\int \varphi^ {*} H \varphi d \tau \geqslant 0,\tag{4.48}
\end{equation*}

则算子就称为正的。如果算子 $H$ 是埃尔米特的，则 $\varphi$ 应属于容许函数类。例如，算子 $H = -\triangle$ ，当它所作用的函数满足边界条件 $\varphi = 0$ 时，它是正的。因为在这种情形

\begin{equation*}
- \int \varphi^ {*} \triangle \varphi d \tau = \int | \nabla \varphi | ^ {2} d \tau \geqslant 0.
\end{equation*}

后面的关系式是通过分部积分并利用边界条件而得到的。类似地，如果 $f \geqslant 0$ ，则算子 $H = -\nabla [f\nabla ]$ 是正的。因为此时

\begin{equation*}
- \int \varphi^ {*} \nabla [ f \nabla ] \varphi d \tau = \int f | \nabla \varphi | ^ {2} d \tau \geqslant 0.
\end{equation*}

\textbf{定理 4.6}\quad 将只考虑埃尔米特算子。设

\begin{equation*}
H = H _ {0} + H _ {1},\tag{4.49}
\end{equation*}

这里 $H_{1}$ 是正算子；其次设

\begin{equation*}
H _ {0} \varphi_ {n} ^ {0} = \lambda_ {n} ^ {0} \varphi_ {n} ^ {0}, H \varphi_ {n} = \lambda_ {n} \varphi_ {n}.\tag{4.50}
\end{equation*}

则
\begin{equation*}
\lambda_ {n} \geqslant \lambda_ {n} ^ {0}.\tag{4.51}
\end{equation*}

\textbf{证明}\quad 设

\begin{equation*}
H _ {\xi} \equiv H _ {0} + \xi H _ {1},\tag{4.52}
\end{equation*}

则
\begin{equation*}
H _ {\xi} \varphi_ {n} ^ {\xi} = \lambda_ {n} ^ {\xi} \varphi_ {n} ^ {\xi},\tag{4.53}
\end{equation*}

这里
\begin{equation*}
\int \left(\varphi_ {n} ^ {\xi}\right) ^ {*} \varphi_ {n} ^ {\xi} d \tau = 1.
\end{equation*}

对我们有兴趣的是， $\xi$ 从0变到1时 $\lambda_{n}^{\xi}$ 的变化。如果 $\xi$ 改变了量 $d\xi$ ，则

\begin{equation*}
\begin{array}{l} \varphi_ {n} ^ {\xi} \longrightarrow \varphi_ {n} ^ {\xi} + \delta \varphi_ {n} ^ {\xi}, \\ \lambda_ {n} ^ {\xi} \longrightarrow \lambda_ {n} ^ {\xi} + \delta \lambda_ {n} ^ {\xi}. \end{array}\tag{4.54}
\end{equation*}

因为标准化条件仍保持，可写

\begin{equation*}
\begin{array}{r l} \delta \int (\varphi_ {n} ^ {\xi}) ^ {*} \varphi_ {n} ^ {\xi} d \tau & = \int [ \delta (\varphi_ {n} ^ {\xi}) ^ {*} \varphi_ {n} ^ {\xi} + (\varphi_ {n} ^ {\xi}) ^ {*} \delta \varphi_ {n} ^ {\xi} ] d \tau = \\ & = 0. \end{array}\tag{4.55}
\end{equation*}

其次。
\begin{equation*}
\begin{array}{r l} \delta \lambda_ {n} ^ {\xi} = & \int [ \delta (\varphi_ {n} ^ {\xi}) ^ {*} H _ {\xi} \varphi_ {n} ^ {\xi} + (\varphi_ {n} ^ {\xi}) ^ {*} H _ {\xi} \delta \varphi_ {n} ^ {\xi} + \\ & + \delta \xi (\varphi_ {n} ^ {\xi}) ^ {*} H _ {1} \varphi_ {n} ^ {\xi} ] d \tau , \end{array}\tag{4.56}
\end{equation*}

而因为
\begin{equation*}
\begin{array}{r l} \int (\varphi_ {n} ^ {\xi}) ^ {*} H _ {\xi} \delta \varphi_ {n} ^ {\xi} d \tau & = \left[ \int \delta (\varphi_ {n} ^ {\xi}) ^ {*} H _ {\xi} \varphi_ {n} ^ {\xi} d \tau \right] ^ {*} = \\ & = \lambda_ {n} ^ {\xi} \int \delta \varphi_ {n} ^ {\xi} (\varphi_ {n} ^ {\xi}) ^ {*} d \tau , \end{array}
\end{equation*}

故
\begin{equation*}
\begin{array}{r l} \delta \lambda_ {n} ^ {\xi} & = \int [ \lambda_ {n} ^ {\xi} \{(\delta \varphi_ {n} ^ {\xi}) ^ {*} \varphi_ {n} ^ {\xi} + \delta \varphi_ {n} ^ {\xi} (\varphi_ {n} ^ {\xi}) ^ {*} \} + \delta \xi (\varphi_ {n} ^ {\xi}) ^ {*} H _ {1} \varphi_ {n} ^ {\xi} ] d \tau = \\ & = \int \delta \xi [ (\varphi_ {n} ^ {\xi}) ^ {*} H _ {1} \varphi_ {n} ^ {\xi} ] d \tau ; \end{array}
\end{equation*}

因此 $\delta \lambda_{n}^{\xi} / \delta \xi \geqslant 0$ ，这是因为 $H_{1}$ 是正算子。由此显然，$\lambda_{n}^{\xi}$ 随着 $\xi$ 的增加而增加，故当 $\xi = 1$ 时即得 $\lambda_{n} \geqslant \lambda_{n}^{0}$ ，定理证毕。

这样一来，如果边界条件不变，而$H$“增加”一正算子，则特征值不会减小。

例　设
\begin{equation*}
H = - \nabla \cdot [ f \nabla ],
\end{equation*}

这里 $f \geqslant f_0$ 和 $f_0 = a_0 (a_0 > 0)$ 。此外假定在某个曲面 $S$ 上 $\varphi = 0$ 。算子 $H$ 可写成

\begin{equation*}
H = - f _ {0} \triangle - \nabla \cdot [ (f - f _ {0}) \nabla ].
\end{equation*}

我们知道， $H_{0}=-f_{0}\triangle$ 有特征值 $\lambda_{1}^{0}\leqslant\lambda_{2}^{0}\leqslant\cdots$ 和构成完备正交系的特征函数。此外算子

\begin{equation*}
H _ {1} = - \nabla \cdot [ (f - f _ {0}) \nabla ]
\end{equation*}

是正的，所以 $\lambda_{n} \geqslant \lambda_{n}^{0}$ 。因此 $H$ 产生完备正交系。

类似地可考虑算子

\begin{equation*}
H = - \Delta + V,
\end{equation*}

这里 $V \geqslant V_{0} = \mathrm{const} > 0$ 。设 $H_{0} = -\triangle + V_{0}$ ，同时在曲面 $S$ 上 $\varphi = 0$ 和

\begin{equation*}
H _ {1} = V - V _ {0} \geqslant 0.
\end{equation*}

$H_{0}$ 的特征值无界，且 $H_{1} \geqslant 0$ ，所以 $H$ 有构成完备正交系的特征函数。
